% =============================================================================
% Annexes — aide-mémoire, glossaire, questions d'entretien
% =============================================================================
\chapter{Aide-mémoire PostgreSQL}

Tout ce que tu as utilisé dans le projet, sur quelques pages, pour réviser la veille d'un entretien.

\section{Définir les données}
\begin{lstlisting}[style=long]
CREATE TABLE prets (
    pret_id    BIGINT        GENERATED BY DEFAULT AS IDENTITY PRIMARY KEY,
    client_id  BIGINT        NOT NULL REFERENCES clients (client_id) ON DELETE RESTRICT,
    montant    NUMERIC(12,2) NOT NULL CHECK (montant > 0),
    statut     VARCHAR(12)   NOT NULL DEFAULT 'sain' CHECK (statut IN ('sain', 'en_retard', 'en_defaut')),
    CONSTRAINT nom_regle CHECK (...),                       -- contrainte nommée, sur plusieurs colonnes
    UNIQUE (client_id, montant)                             -- unicité composite
);
ALTER TABLE prets ADD CONSTRAINT montant_max CHECK (montant <= 100000);
CREATE INDEX idx_prets_client ON prets (client_id);                       -- clé étrangère : à indexer soi-même
CREATE INDEX idx_impayes ON echeances (pret_id, date_echeance) WHERE montant_paye < montant_du;  -- partiel
DROP TABLE IF EXISTS prets CASCADE;
\end{lstlisting}

\section{Transactions}
\begin{lstlisting}[style=long]
BEGIN;  ...  COMMIT;              -- tout ou rien
BEGIN;  ...  ROLLBACK;            -- tester sans rien modifier
SAVEPOINT avant_test;  ...  ROLLBACK TO avant_test;     -- annuler une partie seulement
SELECT ... FROM prets WHERE pret_id = 61 FOR UPDATE;    -- verrouiller une ligne jusqu'à la fin de la transaction
\end{lstlisting}

\section{Interroger}
\begin{lstlisting}[style=long]
WITH etape1 AS (...), etape2 AS (SELECT ... FROM etape1) SELECT ... FROM etape2;   -- CTE
count(*) FILTER (WHERE statut = 'en_defaut')                -- agrégat conditionnel
coalesce(x, 0)   nullif(x, 0)   greatest(a, b)   least(a, b)
CASE WHEN ... THEN ... WHEN ... THEN ... ELSE ... END
DATE '2026-09-30' - DATE '2026-08-15'                        -- 46 (jours, entier)
date_debut + (k - 1) * INTERVAL '1 month'                    -- toujours depuis la date de départ
date_trunc('month', d)   date_trunc('quarter', d)   extract(year FROM d)
generate_series(1, 36)   generate_series(DATE '2025-01-01', DATE '2025-12-01', INTERVAL '1 month')
percentile_cont(0.5) WITHIN GROUP (ORDER BY revenu)          -- médiane
GROUP BY ROLLUP (stage)                                      -- sous-totaux et total
x IS DISTINCT FROM y                                         -- comparaison qui traite NULL comme une valeur
NOT EXISTS (SELECT 1 FROM ... WHERE ...)                     -- préférer à NOT IN (piège des NULL)
\end{lstlisting}

\section{Fonctions de fenêtre}
\begin{lstlisting}[style=long]
ROW_NUMBER() OVER (PARTITION BY client_id ORDER BY date_demande DESC)      -- 1, 2, 3...
RANK()       OVER (ORDER BY pd DESC)                                       -- 1, 2, 2, 4
DENSE_RANK() OVER (ORDER BY pd DESC)                                       -- 1, 2, 2, 3
NTILE(10)    OVER (ORDER BY pd DESC)                                       -- déciles
LAG(x, 1, 0) OVER (PARTITION BY pret_id ORDER BY num_echeance)             -- ligne précédente (0 au bord)
LEAD(x)      OVER (PARTITION BY pret_id ORDER BY date_arrete)              -- ligne suivante
sum(x) OVER (ORDER BY mois ROWS BETWEEN UNBOUNDED PRECEDING AND CURRENT ROW)  -- cumul
avg(x) OVER (ORDER BY mois ROWS BETWEEN 2 PRECEDING AND CURRENT ROW)       -- moyenne mobile
100 * x / sum(x) OVER ()                                                   -- part du total
sum(sum(x)) OVER ()                                                        -- fenêtre sur un agrégat
... WINDOW w AS (PARTITION BY pret_id ORDER BY num_echeance)               -- fenêtre nommée
-- Filtrer sur une fenêtre : CTE, puis WHERE dans la requête suivante (pas de QUALIFY)
\end{lstlisting}

\section{PL/pgSQL}
\begin{lstlisting}[style=long]
CREATE OR REPLACE FUNCTION f(p_x INTEGER) RETURNS NUMERIC
LANGUAGE plpgsql STABLE AS $$
DECLARE v NUMERIC;
BEGIN
    SELECT ... INTO STRICT v FROM ... WHERE ...;          -- exactement une ligne, sinon erreur
    IF v IS NULL THEN RAISE EXCEPTION 'Message %', p_x USING HINT = '...'; END IF;
    RETURN v;
END $$;

CREATE OR REPLACE FUNCTION fn_trg() RETURNS trigger LANGUAGE plpgsql AS $$
BEGIN
    -- NEW, OLD ; BEFORE : RETURN NEW (ou NULL pour ignorer la ligne) ; AFTER : RETURN NULL
    RETURN NEW;
END $$;
CREATE TRIGGER trg BEFORE UPDATE ON t FOR EACH ROW EXECUTE FUNCTION fn_trg();
CREATE TRIGGER trg2 AFTER UPDATE ON t REFERENCING OLD TABLE AS a NEW TABLE AS n
    FOR EACH STATEMENT EXECUTE FUNCTION fn_trg2();

CREATE OR REPLACE PROCEDURE p(p_id BIGINT DEFAULT NULL) LANGUAGE plpgsql AS $$
DECLARE r RECORD; v_nb INTEGER;
BEGIN
    FOR r IN SELECT ... LOOP ... END LOOP;
    UPDATE ...;  GET DIAGNOSTICS v_nb = ROW_COUNT;
    RAISE NOTICE '% lignes', v_nb;
END $$;
CALL p();          DO $$ BEGIN ... END $$;               -- appel ; bloc anonyme

INSERT INTO t (...) SELECT ... ON CONFLICT (cle) DO UPDATE SET col = EXCLUDED.col;   -- upsert
SELECT pg_notify('canal', 'message');   LISTEN canal;                               -- alertes
GRANT SELECT ON ALL TABLES IN SCHEMA public TO role;   REVOKE ALL ON PROCEDURE p FROM PUBLIC;
\end{lstlisting}

\section{Commandes psql et diagnostic}
\begin{center}\small
\begin{tabularx}{\linewidth}{@{}P{4.5cm} L@{}}
\toprule
\code{psql -d base -f fichier.sql} & exécuter un script (ajouter \code{-v ON_ERROR_STOP=1}) \\
\code{\l}, \code{\c base} & lister les bases, se connecter à une base \\
\code{\dt}, \code{\dv}, \code{\df} & lister les tables, les vues, les fonctions \\
\code{\d nom_table} & structure d'une table : colonnes, index, contraintes, triggers \\
\code{\x}, \code{\timing} & affichage vertical ; durée de chaque requête \\
\code{\copy t FROM 'f.csv' CSV HEADER} & importer un CSV depuis le poste client \\
\code{EXPLAIN ANALYZE requete} & plan d'exécution et temps réels \\
\code{ANALYZE} & mettre à jour les statistiques de l'optimiseur \\
\bottomrule
\end{tabularx}
\end{center}

% =============================================================================
\chapter{Glossaire du risque de crédit}

\begin{description}[leftmargin=0pt, itemsep=0.35em, font=\normalfont\bfseries]
\item[AUC] Aire sous la courbe ROC : probabilité qu'un prêt en défaut ait un score plus risqué qu'un prêt
sain. 0,5 = hasard, 1 = parfait.
\item[Calibration] Qualité d'un modèle dont les probabilités prédites correspondent aux fréquences
observées (une PD de 10~\% $\Rightarrow$ environ 10~\% de défauts).
\item[Capital restant dû (CRD)] Part du capital emprunté pas encore remboursée.
\item[Champion / challenger] Le modèle qui décide en production / les modèles concurrents évalués sur
les mêmes prêts.
\item[Courbe de gains] Part cumulée des défauts captés en refusant les clients les plus risqués, décile
par décile.
\item[CRR] Règlement européen sur les exigences de fonds propres des banques (accords de Bâle). Son
article 178 définit le défaut.
\item[Défaut] Retard de plus de 90 jours sur une obligation significative, ou paiement jugé improbable.
\item[Dotation nette] Variation des provisions sur une période ; elle pèse sur le résultat de la banque
(positive : charge ; négative : reprise).
\item[DPD] \emph{Days Past Due} : nombre de jours depuis l'échéance impayée la plus ancienne.
\item[EAD] \emph{Exposure At Default} : montant dû au moment du défaut ; ici, le capital restant dû.
\item[ECL] \emph{Expected Credit Loss} : perte attendue IFRS~9, $PD \times LGD \times EAD$.
\item[Échéance (mensualité)] Somme due à une date donnée : part de capital + part d'intérêts.
\item[Encours] Total du capital restant dû d'un portefeuille.
\item[Gini] $2 \times AUC - 1$ ; indicateur de discrimination usuel en scoring de crédit.
\item[IFRS 9] Norme comptable qui impose de provisionner les pertes attendues dès l'octroi, en trois stages.
\item[LGD] \emph{Loss Given Default} : part de l'exposition perdue en cas de défaut (45~\% ici).
\item[MOB] \emph{Months On Book} : ancienneté d'un prêt en mois, utilisée dans l'analyse par génération.
\item[PAR$x$] \emph{Portfolio At Risk} : part de l'encours portée par les prêts en retard de plus de $x$ jours.
\item[PD] \emph{Probability of Default} : probabilité de défaut, prédite par le modèle de scoring.
\item[Période probatoire] Délai minimal (au moins trois mois) pendant lequel un prêt sorti d'impayé
reste classé en défaut.
\item[Provision] Montant mis de côté pour couvrir une perte attendue.
\item[Rappel] Part des défauts réels que le modèle aurait refusés.
\item[Roll rate] Taux de passage d'une tranche de retard à une autre d'un mois sur l'autre (matrice de
transition).
\item[Seuil de décision] PD au-delà de laquelle le modèle recommande le refus.
\item[SHAP] Méthode qui décompose une prédiction en contributions de chaque variable.
\item[Stage (IFRS 9)] 1 : sain (PD à 12 mois) ; 2 : risque dégradé, plus de 30 jours de retard (PD sur
la durée restante) ; 3 : défaut (PD = 100~\%).
\item[Taux de couverture] Provisions rapportées à l'exposition.
\item[Taux d'effort] Part du revenu consacrée au remboursement ; ici, le montant emprunté rapporté au
revenu annuel.
\item[Vintage (génération)] Ensemble des prêts accordés sur une même période, suivis dans le temps.
\end{description}

% =============================================================================
\chapter{Questions d'entretien}

Des questions réellement posées en entretien de stage data, avec une réponse courte. Entraîne-toi à
répondre à voix haute, puis à écrire la requête sans regarder.

\section{SQL}

\begin{description}[leftmargin=0pt, itemsep=0.55em]
\item[Le prêt le plus récent de chaque client ?] \code{ROW_NUMBER() OVER (PARTITION BY client_id ORDER
BY date_demande DESC)} dans une CTE, puis \code{WHERE rn = 1}. (PostgreSQL offre aussi
\code{DISTINCT ON (client_id)}.)
\item[Les adresses e-mail en double ?] Un regroupement filtré :
\code{SELECT email FROM clients GROUP BY email} \code{HAVING count(*) > 1}.
\item[COUNT(*), COUNT(col), COUNT(DISTINCT col) ?] Toutes les lignes ; les lignes où \code{col} n'est
pas NULL ; les valeurs distinctes non NULL.
\item[Le piège de NOT IN ?] Si la sous-requête renvoie un seul NULL, \code{x NOT IN (...)} ne renvoie
aucune ligne (la comparaison avec NULL est inconnue). Utilise \code{NOT EXISTS}.
\item[Le piège du LEFT JOIN ?] Une condition sur la table de droite placée dans le \code{WHERE}
élimine les lignes sans correspondance : le \code{LEFT JOIN} redevient un \code{INNER JOIN}. Place-la
dans le \code{ON}.
\item[UNION ou UNION ALL ?] \code{UNION} supprime les doublons (donc trie, plus lent) ; \code{UNION ALL}
garde tout.
\item[Une médiane ?] \code{percentile_cont(0.5) WITHIN GROUP (ORDER BY revenu_annuel)} : 55\,000~€ dans
ta base.
\item[Une requête est lente : que fais-tu ?] \code{EXPLAIN ANALYZE} pour voir le plan ; un index sur les
colonnes de filtre et de jointure ; filtrer avant de joindre ; éviter les fonctions sur les colonnes
indexées ; \code{ANALYZE} si les statistiques sont anciennes.
\end{description}

\section{Modélisation et ingénierie}

\begin{description}[leftmargin=0pt, itemsep=0.55em]
\item[Pourquoi ne pas stocker le retard ?] Il change chaque jour sans événement : stocké, il serait faux
dès le lendemain. On le calcule à partir des dates.
\item[Comment historiser une donnée qui change ?] Ne jamais l'écraser : un journal (comme
\code{journal_statuts}), un registre daté (comme les prédictions), ou une table avec des dates de
validité (« dimension à évolution lente de type 2 » dans un entrepôt de données).
\item[Vue ou vue matérialisée ?] La vue est recalculée à chaque lecture (toujours à jour) ; la vue
matérialisée stocke le résultat (rapide, mais à rafraîchir).
\item[Comment charger un million de lignes ?] Avec \code{COPY} (ou \code{to_sql} avec une méthode
\code{COPY}), dans une transaction, puis \code{ANALYZE}. Jamais un million d'\code{INSERT} séparés.
\item[Comment éviter un double paiement ?] Contraintes (\code{montant_paye <= montant_du}), verrou sur
le prêt (\code{FOR UPDATE}), et en pratique un identifiant de transaction unique pour rendre
l'encaissement idempotent.
\end{description}

\section{Risque de crédit}

\begin{description}[leftmargin=0pt, itemsep=0.55em]
\item[Quand un prêt est-il en défaut ?] Plus de 90 jours de retard sur une obligation significative, ou
paiement improbable (CRR, article 178).
\item[Qu'est-ce que l'ECL ?] La perte attendue IFRS~9, $PD \times LGD \times EAD$, avec une PD à 12 mois
en stage 1, sur la durée restante en stage 2, et 100~\% en stage 3.
\item[Pourquoi une LGD de 45~\% ?] C'est la valeur réglementaire de l'approche « fondation » de Bâle
pour une créance non garantie de premier rang. Une banque en approche avancée l'estime sur ses
propres recouvrements.
\item[À quoi sert une matrice de transition ?] À voir la dynamique du risque : ici, 90~\% des prêts à
61--90 jours passent en défaut le mois suivant, donc le recouvrement doit agir plus tôt.
\item[Pourquoi analyser par génération ?] Pour comparer des prêts au même âge, et détecter qu'une
politique d'octroi s'est dégradée, à condition de n'utiliser que des prêts observés assez longtemps.
\end{description}

\section{Machine Learning et MLOps}

\begin{description}[leftmargin=0pt, itemsep=0.55em]
\item[Comment surveiller un modèle en production ?] Sa discrimination (AUC sur les nouveaux prêts dont
on connaît l'issue), sa calibration (PD prédite contre défaut observé par décile), la dérive des
données (indice de stabilité de la population, PSI) et ses effets métier (taux de refus).
\item[Pourquoi l'échantillon de test seulement ?] Sur l'échantillon d'apprentissage, le modèle a déjà vu
les réponses : ses performances y sont surestimées.
\item[À quoi sert SHAP dans le crédit ?] À expliquer chaque décision (obligation envers le client et le
régulateur), à vérifier que le modèle s'appuie sur des variables légitimes, et à détecter des biais.
\item[Pourquoi un XGBoost « monotone » ?] Des contraintes de monotonie garantissent, par exemple, que le
risque prédit ne baisse jamais quand le taux d'effort augmente : le modèle reste cohérent avec
l'intuition métier, donc acceptable par la validation et le régulateur.
\item[Quand remplacer le champion ?] Quand un challenger fait mieux sur les mêmes prêts, sur des données
récentes et hors apprentissage, à risque égal ou plus faible, et que l'écart est stable dans le temps.
Le registre des prédictions permet exactement cette comparaison.
\end{description}
